%% kapa: exam=1 round=33 year=2022 period=1 form=A subject=civil_law subject_ko=민법 seq=7 number=7
%% source: Q-Net 감정평가사 기출문제 (cst00309) artlSeq=5209997 file=raw/1008/5209997/2022년도 제33회 감정평가사 1차시험 1교시(민법,경제학,부동산학).pdf
%% conversion: kordoc 4.24.51 -> tools/gen1.py (auto); review=proofread-sonnet; flags=none
\begin{question}{7}
甲사단법인이 3인의 이사(乙, 丙, 丁)를 두고 있는 경우에 관한 설명으로 옳지 않은 것은? (다툼이 있으면 판례에 따름)
\begin{choices}
  \item 乙, 丙, 丁은 甲의 사무에 관하여 원칙적으로 각자 甲을 대표한다.
  \item 甲의 대내적 사무집행은 정관에 다른 규정이 없으면 乙, 丙, 丁의 과반수로써 결정한다.
  \item 甲의 정관에 乙의 대표권 제한에 관한 규정이 있더라도 이를 등기하지 않으면 그와 같은 정관의 규정에 대해 악의인 제3자에 대해서도 대항할 수 없다.
  \item 丙이 제3자에게 甲의 제반 사무를 포괄 위임한 경우, 그에 따른 제3자의 사무대행행위는 원칙적으로 甲에게 효력이 없다.
  \item 甲의 토지를 丁이 매수하기로 한 경우, 이 사항에 관하여 丁은 대표권이 없으므로 법원은 이해관계인이나 검사의 청구에 의하여 임시이사를 선임하여야 한다.
\end{choices}
\end{question}
\answer{5}
